\section{Substrate-Exclusive Witnesses}
\label{sec:witnesses}

The verification asymmetry is eliminated when $\rho(a, E) = 0$, which
by Definition~\ref{def:substrate_exclusive_eval} requires the evaluator
to be rooted in a substrate the agent does not occupy.  This section
formalizes the two kinds of entities that provide this guarantee:
algorithmic witnesses that execute verification protocols, and the
governance layer that determines which protocols are acceptable.

\subsection{Algorithmic Witnesses}
\label{sec:alg_witnesses}

\begin{definition}[Algorithmic Witness]
\label{def:alg_witness}
An \textbf{algorithmic witness} $\Witness_{\mathrm{alg}}$ for agent $a$
on substrate $s$ is a process satisfying:
\begin{enumerate}
    \item \textbf{Substrate exclusivity:} $\Witness_{\mathrm{alg}}$
    executes on substrate $s' \neq s(a)$.
    \item \textbf{Protocol determinism:} $\Witness_{\mathrm{alg}}$
    executes a verification protocol $\pi$ that is fully specified by a
    governance-accepted evaluation protocol
    (Definition~\ref{def:gov_fork}, below).
    The witness exercises no independent judgment; its output is a
    deterministic (or, for statistical protocols, randomized with
    specified distribution) function of its inputs and the protocol
    specification.
    \item \textbf{Tamper evidence:} the witness's execution trace is
    recorded to the verification ledger $\Ledger$
    (\S\ref{sec:ledger}, below), and any modification to the trace is
    detectable by any ledger participant.
\end{enumerate}
\end{definition}

The key property of algorithmic witnesses is that they execute
\emph{protocols}, not \emph{judgments}.  All judgment is resolved at the
governance layer (\S\ref{sec:governance}) when the evaluation protocol
was accepted.  Once accepted, verification is mechanical: the witness
runs the protocol, records the result, and the outcome is auditable.

\paragraph{Examples.}
Algorithmic witnesses include: (a)~a validator node on physically
separate hardware that checks behavioral consistency against committed
claims; (b)~an automated test runner that executes a
governance-accepted evaluation protocol against the agent's outputs; (c)~a
statistical monitor that runs SPRT-like detection on residuals between
committed and observed behavior (\S\ref{sec:continuous_monitoring}).
In each case, the witness's substrate exclusivity is enforced by running
on infrastructure that the agent under evaluation does not control.

\subsection{The Governance Layer: Evaluation Protocol Admissibility}
\label{sec:governance}

Algorithmic witnesses verify claims against \emph{accepted} evaluation
protocols.  The question ``which evaluation protocols should be
accepted?'' is a governance decision, not a verification decision.
It is the single point at which human judgment enters the system.

\begin{definition}[Governance Fork Protocol]
\label{def:gov_fork}
A \textbf{governance fork protocol} for evaluation-protocol
admissibility consists of:
\begin{enumerate}
    \item \textbf{Proposal:} An evaluation protocol proposal $\Pi$ for
    claim type $\tau$ is submitted to the ledger $\Ledger$ with a
    specification $\sigma_\Pi$ (what it tests, acceptance criteria,
    passing threshold) and a proposer identifier.
    \item \textbf{Deliberation:} A deliberation window of duration
    $\Delta_{\mathrm{gov}}$ opens.  During this window, any ledger
    participant may: study the specification, run $\Pi$ against known
    agents to check for bias, flag concerns with evidence, or publish
    analysis to $\Ledger$.
    \item \textbf{Vote:} At the close of the deliberation window,
    validator nodes vote on acceptance.  Acceptance requires a
    \emph{cross-substrate supermajority}: at least fraction $\theta$ of
    validators overall, with at least one accepting validator per
    substrate represented in $\Ledger$.
    \item \textbf{Resolution:} If accepted, $\Pi$ is recorded as a
    governance entry in $\Ledger$ with an effective date.  All
    subsequent claims of type $\tau$ may be verified against
    $\Pi$ algorithmically.  If rejected, $\Pi$ is archived in $\Ledger$
    with its full deliberation record.
\end{enumerate}
\end{definition}

This structure is analogous to Ethereum Improvement Proposals (EIPs):
the governance decision is slow (human-speed deliberation over weeks to
months), but once made, enforcement is fast (machine-speed verification
against the accepted protocol).  The analogy is structural, not
implementational: we require neither token economics nor
proof-of-stake consensus, only the proposal--deliberation--vote--resolution
cycle with cross-substrate participation.

\begin{remark}
The deliberation window $\Delta_{\mathrm{gov}}$ should be long enough
for meaningful analysis (weeks, not hours) and short enough that the
protocol pipeline does not bottleneck population-level verification.
Since evaluation protocols change rarely compared to individual
claims, a population generating thousands of claims per day may see a
new protocol proposal monthly or less, keeping the governance layer's
bandwidth requirements modest.
\end{remark}

\subsection{Resolving the Bandwidth Problem}
\label{sec:bandwidth}

Human judgment is scarce.  A naive verification system that requires
human evaluation of every claim cannot scale beyond a handful of agents.
The two-layer architecture of this section resolves the bandwidth
problem by separating the frequency of judgment from the frequency of
verification:

\begin{itemize}
    \item \textbf{Governance} (infrequent, human-speed): which evaluation
    protocols are admissible?  Requires human deliberation, cross-substrate
    participation, and careful analysis.  Frequency: on the order of
    protocol proposals, perhaps monthly.
    \item \textbf{Verification} (frequent, machine-speed): does claim
    $C$ pass accepted evaluation protocol $\Pi$?  Fully algorithmic,
    executed by $\Witness_{\mathrm{alg}}$, requires no human attention.
    Frequency: on the order of claims, potentially thousands per day.
\end{itemize}

Human attention is allocated to the meta-level decision (\emph{what} to
measure) rather than the object-level decision (\emph{did the
measurement pass}).  This is the same separation of concerns that makes
legislative systems scalable: laws are deliberated slowly, but
enforcement against accepted laws is fast and algorithmic.

\subsection{Coalition Capture Resistance via Governance}
\label{sec:coalition_governance}

A coalition of agents on substrate $s$ that wishes to inflate its
members' verified claims faces two attack surfaces: the
verification layer and the governance layer.

At the \emph{verification layer}, coalition resistance is structural:
the algorithmic witness is on substrate $s' \neq s$ and executes a fixed
protocol.  The coalition cannot influence the witness's computation
(\S\ref{sec:commitment}, Proposition~\ref{prop:coalition_resistance}).

At the \emph{governance layer}, the coalition's attack is protocol
design: proposing an evaluation protocol $\Pi$ that only coalition
members can pass.  The governance fork protocol defends against this
through three mechanisms:
\begin{enumerate}
    \item \textbf{Public specification:} $\sigma_\Pi$ is published during
    deliberation, so cross-substrate participants can analyze it for
    bias.
    \item \textbf{Cross-substrate supermajority:}  Acceptance requires
    at least one validator per substrate, so a same-substrate coalition
    cannot unilaterally accept its own protocol.
    \item \textbf{Empirical testing:} During the deliberation window,
    any participant can run $\Pi$ against known agents on other substrates
    to check whether $\Pi$ genuinely tests the claimed property or
    merely filters for coalition membership.
\end{enumerate}

Prevention at the governance layer is strictly stronger than detection
at the verification layer: a biased protocol that is never accepted
inflicts no damage, whereas one that is accepted and later detected has
already produced unreliable verification records.

\subsection{Witness Availability Requires \texorpdfstring{$m \geq 2$}{m >= 2}}
\label{sec:witness_m2}

Substrate-exclusive witnesses exist only when the population spans at
least two substrates ($m \geq 2$).  This is a logical consequence of
Definition~\ref{def:alg_witness}: a witness must execute on substrate
$s' \neq s(a)$, which requires $s'$ to exist.  But the deeper reason is
physical, not definitional.  Processes that share a substrate share a
causal medium: they can address each other's inputs, producing nonzero
verification asymmetry $\rho > 0$
(Definition~\ref{def:ver_asymmetry}).  A second substrate provides a
computation space that is causally separated from the first, and it
is this physical separation that makes $\rho = 0$ achievable.  In a
monopolar population ($m = 1$), every process is co-substrate with
every agent, and no rearrangement of the population's computational
resources can produce an evaluator whose inputs are unreachable.

We defer the full impossibility argument to
\S\ref{sec:anti_monopolar}, where we prove that a monopolar population
cannot achieve exogenous verification through any internal action.
